\begin{textsample}{NEG Dim 9 – global\_south\_2025\_04 – Score -3.00 – global\_south\_2025\_04/123230684.txt}  \label{ex:f9_neg_266}
Mobile phone footage has emerged that appears to \textbf{contradict} Israel ' s account of why soldiers opened fire on a convoy of ambulances and a fire truck on March 23, killing 15 rescue workers.
The video, published by the Palestine Red Crescent Society ( PRCS ), shows the vehicles moving in darkness with headlights and emergency flashing lights switched on - before coming under fire.
The PRCS said the video was obtained from the phone of a paramedic who was killed.
The Israel Defense Forces ( IDF ) initially denied the vehicles had their headlights or emergency signals on.
But in response to the new video, the IDF told the BBC :"All claims, including the documentation circulating about the incident, will be thoroughly and deeply examined to understand the sequence of events and the handling of the situation".
A surviving paramedic previously told the BBC that the ambulances were clearly marked and had their internal and external lights on.
The latest video, which the PRCS said had been shown @ @ @ @ @ @ @ @ @ @ drawing to a halt on the edge of the road, lights still flashing, and at least two emergency workers stepping out wearing reflective clothing.
The windscreen of the vehicle being filmed from is cracked and shooting can then be heard lasting for several minutes as the person filming says prayers.
He is understood to be one of the dead paramedics.
The footage was found on his phone after his body was recovered from a shallow grave one week after the incident.
The bodies of the eight paramedics, six Gaza Civil Defence workers and one UN employee were found buried in sand, along with their wrecked vehicles.
It took international organisations days to negotiate safe access to the site.
Israel claimed a number of Hamas and Islamic Jihad militants had been killed in the incident, but it has not provided any evidence or further explained the threat to its troops.
The IDF promised to investigate the circumstances after a surviving paramedic questioned its account.
In an interview with the BBC, paramedic Munther @ @ @ @ @ @ @ @ @ @ it ' s the same thing.
External and internal lights are on.
Everything tells you it ' s an ambulance vehicle that belongs to the Palestinian Red Crescent.
All lights were on until the vehicle came under direct fire."He also denied he or his team had any militant connections."All crews are civilian.
We don't belong to any militant group.
Our main duty is to offer ambulance services and save people ' s lives.
No more, no less,"he said.
Speaking at the United Nations yesterday the President of the PRCS, Dr Younis Al-Khatib, referred to the video recording, saying :"I heard the voice of one of those team members who was killed.
His last words before being shot... ' forgive me mum, I just wanted to help people.
I wanted to save lives '.
It ' s heartbreaking".
He called for"accountability"and"an"independent and thorough investigation"of what he called an"atrocious crime @ @ @ @ @ @ @ @ @ @ the 23 March incident.
DISCLAIMER : The Views, Comments, Opinions, Contributions and Statements made by \textbf{Readers} and Contributors on this platform do not \textbf{necessarily} represent the views or policy of Multimedia Group Limited.

% matched lemmas: contradict, necessarily, reader
\end{textsample}
